--codigo de Latex--

En la segunda parte del trabajo practico se analizaron dos configuraciones de un circuito de corriente continua con dos mallas y dos fuentes, con el objetivo de verificar experimentalmente las Leyes de Kirchoff y comparar los valores medidos con los obtenidos teoricamente

Para la primera configuracion, con la fuente B2 (9V) en su orientacion original, se observo que las corrientes medidas en las ramas asociadas a B1 y R1 fueron de 35,8mA, mientras que el valor teorico obtenido fue de 38,54mA. En la rama central, con la resistencia R2, se midio una corriente de 26,2mA, valor muy cercano al teorico de 25,955mA. Finalmente, en la rama asociada a R3 y B2 se registro una corriente de 9,58mA, mientras que la teorica de 12,588mA. Por lo tanto, en la primer parte se obtuvo una corriente aceptable, aunque con diferencias mas notorias en la rama de menor corriente.

Para la segunda configuracion, al invertir la fuente B2 (9V), se observo una concordancia general mas cercana entre los resultados experimentales y los teoricos. En la rama de B1 y R1 se midio una corriente de 42,7mA, frente a un valor teorico de 43,947mA. En la rama central, correspondiente a R2, se obtuvo 23mA, con un valor teorico de 23,475mA. Finalmente, en la rama de R3 y B2 se midio 19,7mA, tambien cercano al valor calculado de 20,468mA. En este segundo experimento, se pudo verificar las Leyes de Kirchoff con una exactitud levemente mayor a la primera configuracion.

En ambos casos, se comprobo cualitativamente la ley de los nodos, ya que la corriente de la rama central resulto coherente con la diferencia entre las ramas asociadas a R1 y R3, de acuerdo con la conservacion de la carga. La comparacion entre los valores teoricos y experimentales mostro que el circuito respondio de manera consistente con la ley de malla, es decur, con la conservacion de la energia en cada recorrido cerrado.

La diferencias observadas entre los valores medidos y calculados pudieron atribuirse a diversas causas experimentales, entre ellas la tolerancia de las resistencias utilizadas, la resistencia interna de las fuentes, la precision finita de los multimetros y posibles errores de lectura. Ademas, al cambiar la escala de uno de los amperímetros se observó una variación en las lecturas, lo cual pudo explicarse considerando que los instrumentos no son ideales y que su resistencia interna depende de la escala seleccionada, alterando levemente las condiciones del circuito.

En síntesis, los resultados obtenidos permitieron verificar las leyes de Kirchhoff y comparar experimentalmente la distribución de corrientes en ambas configuraciones del circuito. Se observó, además, que la inversión de una de las fuentes modificó la circulación de corriente tal como se esperaba teóricamente.
